\documentclass[letterpaper,10pt]{article}

% Generated from site/src/data/cv.js and profile.js by "npm run cv-pdf".
% Do not edit this file: change the data and run the command again.

\usepackage[left=0.48in,top=0.42in,right=0.48in,bottom=0.42in]{geometry}
\usepackage{enumitem}
\usepackage{parskip}
\usepackage{hyperref}
\usepackage[T1]{fontenc}
\usepackage{lmodern}
\usepackage{titlesec}

\hypersetup{hidelinks}

\setlength{\parindent}{0pt}
\setlength{\parskip}{0pt}

\titleformat{\section}
    {\large\bfseries}
    {}
    {0em}
    {}
    [\vspace{-5pt}\hrule\vspace{3pt}]

\titlespacing*{\section}{0pt}{7pt}{3pt}

\setlist[itemize]{
    leftmargin=0.16in,
    itemsep=1.2pt,
    topsep=1pt,
    parsep=0pt,
    partopsep=0pt
}

\newcommand{\resumeItem}[1]{
    \item \small #1
}

\begin{document}

\begin{center}
    {\LARGE \textbf{Gautham Sam}} \\[3pt]

    \small
    Hyderabad, India \quad | \quad
    +91 80882 29827 \quad | \quad
    \href{mailto:gauthammsam57@gmail.com}{gauthammsam57@gmail.com} \\

    \href{https://www.linkedin.com/in/gautham-sam/}{linkedin.com/in/gautham-sam} \quad | \quad
    \href{https://github.com/gauthammsam57}{github.com/gauthammsam57}
\end{center}

\section*{PROFESSIONAL SUMMARY}

\small
I work on the software that makes robots useful: mobile robots, arms, the arm of a humanoid, and the devices and platforms around them. I trained as a mechanical engineer and have spent six years in robotics, first on care robots at Invento Robotics, now at Neural Foundry, where I was the first robotics engineer and work with a team of six.

\section*{AT A GLANCE}

\small
\textbf{Years in robotics:} 6 (since 2020)

\textbf{Robot types:} Mobile robots, collaborative arms, humanoid arm, lab devices

\textbf{Current team:} 6 engineers

\textbf{Pull requests:} 212 written, 345 reviewed for teammates

\textbf{Main stack:} ROS 2, Python, C++, embedded C

\section*{WHAT I WORK ON}

\small
\textbf{Robotics architecture} -- Choosing robot platforms, system architecture, joining hardware and software, reusable platforms, several robots at once

\textbf{Technical direction} -- Engineering standards, architecture and code reviews, technical specifications, mentoring, hiring, release practice

\textbf{Automation} -- Lab automation, device orchestration, automation languages, workflow execution, recovery and fault handling

\textbf{Systems} -- Embedded systems, ROS 2, distributed systems, simulation, teleoperation, cloud infrastructure, deployment

\section*{PROFESSIONAL EXPERIENCE}

\textbf{Neural Foundry}
\hfill \textbf{June 2024 -- present}\\
\textit{Robotics Lead} (joined as Senior Robotics Engineer)
\hfill Hyderabad, India

\begin{itemize}
    \resumeItem{First robotics engineer at the company. The team is now six; I have been part of hiring each person, and built our take-home test, with its own simulator and scoring.}
    \resumeItem{First engineer on ALEX (Automated Lab Executor), a ROS 2 system that runs lab tasks with a collaborative robot arm. 482 commits to it in its first nine months; first stable release closed in early 2026.}
    \resumeItem{ALEX now has two generations of its automation language, workflow execution, device actions, telemetry, cancel, pause and resume, and it can pick a long lab protocol back up after a failure.}
    \resumeItem{With the team, brought around ten lab devices under robot control in a client's lab: syringe pumps, gravimetric liquid and solid dispensers, balances, a pipette station, a homogenizer, a sonicator. Each follows the same path from bare hardware to a scriptable action: firmware, control node, software abstraction, device management, workflow.}
    \resumeItem{Designed the ROS 2 client and the embedded client that connect robots and devices to URP, the company's robotics platform: a config-driven command server for arms, mobile robots and cameras; a telemetry and command registry; device logs as telemetry; authenticated MQTT over TLS. A new robot is brought on by configuration, not by a build of its own.}
    \resumeItem{Found the same control server written three times, for arms, mobile robots and audio-video, and merged them into one shared base without changing what they do.}
    \resumeItem{Designed a hardware abstraction for the embedded code, so the same application logic runs on a workstation simulator, an Arduino and an ESP32, and firmware can be tested without the hardware.}
    \resumeItem{Made it possible to run several robots at once: each robot's stack is isolated and has a virtual mode, so many can work against one platform and one simulation.}
    \resumeItem{Built safe-path planning for the arm (a multi-layer graph of safety points, searched with Dijkstra's algorithm) and fault recovery (runtime errors, reconnection, resuming a task partway).}
    \resumeItem{Designed the peer-to-peer link for RAVEN, the company's VR teleoperation of a robot arm. Tested Europe to India over a home connection at about 100 ms.}
    \resumeItem{Moved robot camera streaming from a custom WebRTC client to WHIP with a media server, so one published stream serves every viewer; added stereo camera support.}
    \resumeItem{Ran the move of the whole system to ROS 2 Jazzy: build environment, robot application, Doosan driver, development environment and automation language, across more than 200 files.}
    \resumeItem{Look after the project building the system's simulated twin in Isaac Sim, and have reviewed 65 of its pull requests.}
    \resumeItem{Plan two-week sprints, review code, write architecture documents, run releases from development to production.}
    \resumeItem{Support client work: site visits, robot selection (arms, SCARA, mobile robots from several vendors), audit responses, proposals.}
\end{itemize}

\vspace{2pt}

\textbf{Invento Robotics (Mitra Robot)}
\hfill \textbf{January 2023 -- May 2024}\\
\textit{Senior Software Engineer}
\hfill Bangalore, India

\begin{itemize}
    \resumeItem{Built joystick teleoperation for the robot fleet over ROS Bridge, with vibration feedback when the robot nears an obstacle.}
    \resumeItem{Built a Python WebSocket server and client for robot-to-portal communication, tested to over 1,000 concurrent connections.}
    \resumeItem{Worked with the chief architect on the image and video pipelines for Mitra AI (FastAPI, Celery with Redis, Milvus vector search).}
    \resumeItem{Mapped and deployed robots at customer sites, including mapping one site remotely from another country.}
    \resumeItem{In the care portal, did most of the work on the live map after the founder started it: robot position and lidar on a canvas, sending navigation points, a path layer, saved patrols.}
    \resumeItem{Wrote the first version of a React Native interface app for the robot: speech, camera, video and option screens driven by defined flows.}
    \resumeItem{Set up a relay server so robots on site networks could be reached remotely, and a logging service (FastAPI, Elasticsearch).}
    \resumeItem{With the team, built a person-following and gesture-response behaviour, and a personalised daily flow of news, music and exercise for seniors.}
    \resumeItem{Built recording and cleaning tools for BodyLangRec, a public gesture dataset (35 gestures, 39 participants, 22 of them over 70).}
    \resumeItem{Moved server infrastructure from Google Cloud to AWS with no downtime.}
    \resumeItem{Built bots for a machine learning course: question answering, quizzes, onboarding.}
\end{itemize}

\vspace{2pt}

\textbf{Invento Robotics (Mitra Robot)}
\hfill \textbf{October 2020 -- December 2022}\\
\textit{Software Engineer}
\hfill Bangalore, India

\begin{itemize}
    \resumeItem{With the team, rebuilt an inherited ROS 1 navigation stack and migrated it to ROS 2 and Nav2. Wrote path-tracking controllers (pure pursuit, Stanley). The robot later ran 1,000 km autonomously without an incident, as the company announced.}
    \resumeItem{Wrote autonomous docking using ArUco markers; trained teammates and customers on it.}
    \resumeItem{Added a fleet patrol feature.}
    \resumeItem{Worked on the takeover page of the fleet web app: driving by gamepad and keyboard, chat, robot controls.}
    \resumeItem{Prototyped the control system for a humanoid arm and hands: Arduino motor control, I2C, power measurement.}
    \resumeItem{Worked on the robot's iOS app in Swift: the socket link to the fleet, Wi-Fi setup over Bluetooth, text-to-speech voices with caching, VoIP calling.}
    \resumeItem{Built tools for preparing the robot's conversation content.}
\end{itemize}

\vspace{2pt}

\textbf{Invento Robotics (Mitra Robot)}
\hfill \textbf{August -- September 2020}\\
\textit{Software Engineer Intern}
\hfill Bangalore, India

\begin{itemize}
    \resumeItem{Robot URDF model and Gazebo simulation; test procedures for the navigation stack; first draft of the technical manual.}
\end{itemize}

\vspace{2pt}

\textbf{Earlier internships}

\begin{itemize}
    \resumeItem{Spatez Technology, Thrissur \textperiodcentered{} December 2019 -- January 2020}
    \resumeItem{Innovians Technologies, Noida \textperiodcentered{} June -- August 2019}
\end{itemize}

\section*{TECHNICAL SKILLS}

\small
\textbf{Languages:} Python, C++, embedded C, Bash, JavaScript, Swift

\textbf{Robotics:} ROS 2 (Humble, Jazzy), Nav2, MoveIt 2, rosbridge, URDF

\textbf{Simulation:} Isaac Sim, Gazebo

\textbf{Robots and hardware:} Doosan collaborative arms, mobile robot bases, ZED and RealSense cameras, ESP32, Arduino

\textbf{Device protocols:} RS485, RS232, Modbus, UART, I2C, TCP/IP, RTSP

\textbf{Messaging and streaming:} MQTT, WebSocket, WebRTC, WHIP

\textbf{Backend and cloud:} FastAPI, Celery, Redis, Elasticsearch, Docker, AWS, Google Cloud

\textbf{AI and vision:} OpenCV, tools built on language models, vector search (Milvus)

\textbf{Team practice:} Sprint planning, code review, architecture documents, hiring

\section*{OPEN SOURCE}

\small
MoveIt 2: fix to the simulated Panda servo configuration (merged, 2025).

Doosan \texttt{doosan-robot2}: variable fix (merged, 2026).

\section*{PUBLICATIONS AND DATA}

\small
``Low Cost FGC Electric Wheelchair'', RSP Science Hub: a gesture-controlled wheelchair built at college.

BodyLangRec: A Gesture Dataset (Kaggle), contributor.

\section*{COURSES AND PRIZES}

\small
NPTEL Robotics, Elite Silver: kinematics, dynamics, motion and trajectory planning.

Paper presentations: first prize at CSI College of Engineering, second prize at Vellalar College of Engineering.

\section*{COMMUNITY}

\small
ROSCon India 2025, Pune.

\section*{EDUCATION}

\small
B.E. Mechanical Engineering, CSI College of Engineering (Anna University), 2016 -- 2020. CGPA 6.7 / 10.

\end{document}
